The empirical results generated by the simulation provided definitive answers to the three research questions posed at the outset of this study:

\noindent{\textbf{RQ1: Impact of Behavioral Parameters.}}\\
The global sensitivity analysis and comparative scenario testing confirmed that variations in household behavioral parameters fundamentally dictated policy stability. Specifically, the simulation revealed that the ``Cost of Effort'' was the dominant inhibitor of compliance, overriding high ``Attitude'' scores \citep{Zheng2022}. This explained why previous IEC campaigns failed; increasing awareness (Attitude) did not compensate for the high logistical friction (PBC) in areas like Poblacion \citep{Moeini2023}. Consequently, policies that did not physically reduce this cost—either through convenient infrastructure or high-intensity enforcement that made non-compliance ``costlier''—were destined to fail, regardless of public awareness levels \citep{Chen2023, Badua2022}.

\noindent{\textbf{RQ2: Optimal Resource Allocation Ratio.}}\\
The study determined that there was no single static ``golden ratio'' for resource allocation. Instead, the optimal allocation was \textit{dynamic and state-dependent}. The HuDRL agent demonstrated that the most effective strategy involved a phased approach: an initial heavy concentration (up to 90\%) on Enforcement and Incentives to break resistance, followed by a gradual pivot toward IEC and maintenance spending once the ``tipping point'' of 70\% compliance was reached \citep{Nishimura2022}. This challenged the static 20\%/30\%/50\% splits often found in annual appropriation ordinances, proving that allocation ratios had to evolve quarter-by-quarter based on real-time compliance data \citep{Abdallah2021, Torkayesh2021}.

\noindent{\textbf{RQ3: The Superior Dynamic Strategy.}}\\
The Heuristic-Guided PPO algorithm was unequivocally identified as the dynamic strategy yielding the highest cost-benefit ratio. While the ``Status Quo'' achieved only 12.5\% global compliance and ``Pure Enforcement'' collapsed to near-zero in key areas, the DRL-driven ``Sequential Saturation'' strategy achieved 100\% graduation of all barangays within 12 quarters without exceeding the budget cap \citep{Jimenez2025}. This confirmed that the highest compliance per peso was achieved not by spending more, but by spending sequentially—concentrating capital to trigger social norms that subsequently sustained behavior for free \citep{ZhengAI2022, Andre2021}.